% NOTE: This template uses IEEEtran (ICDM format).
% For KDD format instead, swap to ACM acmart:
%   \documentclass[sigconf]{acmart}
% and replace \IEEEauthorblockN/\IEEEauthorblockA with acmart's \author/\affiliation.
\documentclass[conference]{IEEEtran}
\IEEEoverridecommandlockouts
\usepackage{cite}
\usepackage{amsmath,amssymb,amsfonts}
\usepackage{algorithmic}
\usepackage{graphicx}
\usepackage{textcomp}
\usepackage{xcolor}
\usepackage{CJKutf8}
\usepackage{multirow}

\begin{document}

\title{Graph Neural Networks for Fraud Detection: A Survey and an
Imbalance-Aware Approach}

\author{\IEEEauthorblockN{Zamo Rzgar Ahmed}
\IEEEauthorblockA{\textit{Department of Software Engineering} \\
\textit{Beihang University (BUAA)}\\
Beijing, China \\
lb2606201@buaa.edu.cn}
\and
\IEEEauthorblockN{Course Instructor: \begin{CJK}{UTF8}{gbsn}张日崇\end{CJK} (Richong Zhang)}
\IEEEauthorblockA{\textit{Course: Data Mining} \\
\textit{Beihang University (BUAA)}\\
Beijing, China}
}

\maketitle

\begin{abstract}
Fraud detection in online services can be formulated as semi-supervised node
classification on graphs, where fraudsters are heavily outnumbered and
actively camouflage themselves among benign users. This report surveys
representative graph-based fraud detection methods --- from the dense-block
detector FRAUDAR to graph neural network (GNN) approaches such as SemiGNN,
CARE-GNN, and PC-GNN --- and proposes an imbalance-aware GNN that combines
label-supervised similarity filtering of camouflaged neighbors with balanced
neighborhood sampling. Experiments on the YelpChi and Amazon fraud benchmarks
show that the proposed model outperforms a GCN baseline by a large margin and
surpasses SemiGNN on YelpChi across all metrics (0.854 vs.\ 0.801 AUROC,
mean over three seeds), while remaining competitive on Amazon. Ablation and
interaction studies show that the similarity filter drives the improvement,
whereas balanced sampling adds no measurable gain once class-weighted
training is used. Code and results are publicly available.
\end{abstract}

\begin{IEEEkeywords}
graph neural networks, fraud detection, anomaly detection, class imbalance,
semi-supervised learning
\end{IEEEkeywords}

\section{Introduction}
Online review and e-commerce platforms lose revenue and user trust to
fraudulent accounts. Because fraudsters deliberately mimic benign behavior
(\emph{camouflage})~\cite{DBLP:conf/kdd/HooiSBSSF16}, purely feature-based
detectors are easily evaded. Modeling users, reviews, and products as a graph
and applying graph neural networks (GNNs) has therefore become a dominant
paradigm; deep graph anomaly detection more broadly is surveyed
in~\cite{DBLP:journals/tkde/MaWXYZSXA23}. Two challenges are central: (1)
extreme class imbalance between fraudsters and benign users, and (2)
camouflage along both feature and relation dimensions. This project studies
GNN-based fraud detectors that explicitly address these challenges.

Concretely, this report makes three contributions. First, it surveys
representative graph-based fraud detection methods from top venues and
organizes them along the two challenges above (Section~II). Second, it
proposes a simple imbalance-aware GNN that replaces CARE-GNN's
reinforcement-learning neighbor selector with a fixed, label-supervised
similarity filter, and combines it with PC-GNN-style balanced sampling
(Section~IV). Third, it evaluates the model on two public benchmarks with
three random seeds and reports ablation and interaction studies that isolate
the contribution of each component (Section~V).

\section{Related Work}
\subsection{Graph-based Fraud Detection without Deep Learning}
A long line of work detects fraud through purely structural signals:
groups of fraudulent accounts act in lockstep, creating unusually dense
blocks in the adjacency matrix of, e.g., a user--product review graph.
Density-based detectors exploit exactly this signal, but fraudsters
evade them through \emph{camouflage}: adding edges to honest, popular
targets so as to ``look normal.'' FRAUDAR~\cite{DBLP:conf/kdd/HooiSBSSF16}
is the representative method of this family. It detects dense subgraphs
with a greedy, near-linear-time algorithm whose suspiciousness metric is
designed to be robust to camouflage edges, and it provides provable,
data-dependent bounds on the amount of fraud that can remain
undetected on a given graph. In experiments, FRAUDAR detected injected
fraud on a Twitter graph with high accuracy even under camouflage, where
earlier density-based methods fail. Such methods are scalable and effective,
but they ignore node features and the (scarce) available labels, which
limits their accuracy on attributed graphs.

\subsection{Deep Graph Anomaly Detection}
Ma et al.~\cite{DBLP:journals/tkde/MaWXYZSXA23} systematize deep learning
approaches to graph anomaly detection, organizing methods by the
anomalous object they target (node, edge, subgraph, or graph level) and
by architecture. Their survey highlights two challenges that are central
to fraud detection and to this work: (1) extreme \emph{class imbalance},
since anomalies are by definition rare, and (2) \emph{camouflage}, since
adversarial anomalies deliberately resemble normal instances. Fraud
detection on review or transaction networks corresponds to node-level
anomaly detection on attributed, multi-relation graphs under
semi-supervision --- the setting adopted by all methods below.

\subsection{GNN-based Fraud Detectors}
Graph neural networks (GNNs) detect fraudulent nodes by aggregating
neighborhood information. SemiGNN~\cite{DBLP:conf/icdm/WangQL0JWFYZY19}
is a semi-supervised graph attentive network developed for
financial fraud detection at payment-platform scale: it organizes each
node's neighborhood into multiple views, applies hierarchical attention
(node-level attention within a view, then view-level attention across
views), and exploits unlabeled nodes through a graph-regularized
semi-supervised loss that pulls together embeddings of adjacent nodes.
This semi-supervision is essential in practice, since labeled fraudsters
are scarce and expensive to obtain.

CARE-GNN~\cite{DBLP:conf/cikm/DouL0DPY20} directly targets camouflage,
distinguishing \emph{feature camouflage} (fraudsters craft feature vectors
similar to benign users) and \emph{relation camouflage} (fraudsters link to
benign entities). It introduces a label-aware similarity measure to identify
informative neighbors, uses reinforcement learning (RL) to adapt a
per-relation filtering threshold, and aggregates the filtered neighborhoods
relation-wise. CARE-GNN also released the YelpChi and Amazon fraud
benchmarks used in this study, on which it clearly outperforms vanilla GCN
and GAT. Its main drawbacks are the cost and instability of the RL-based
selector and that class imbalance is handled only implicitly.

\subsection{Imbalance-aware Graph Learning for Fraud Detection}
Fraudsters typically constitute well under 15\% of nodes, so naive GNN
training is dominated by the majority class. PC-GNN~\cite{DBLP:conf/www/LiuAQCFYH21}
addresses this with a ``pick and choose'' strategy: a label-balanced
sampler \emph{picks} a balanced subset of labeled nodes for each training
step, and a neighborhood sampler \emph{chooses} neighbors by thresholding
feature-space distances per relation, so that aggregation uses a balanced
and informative neighborhood. It reports consistent gains on the YelpChi
and Amazon benchmarks under skewed label distributions. A complementary,
architecture-agnostic line of work handles imbalance at the loss level
with class-weighted or focal objectives; we adopt a class-weighted
cross-entropy in this work and study how it interacts with sampling-level
corrections.

\subsection{Summary and Positioning}
Table~\ref{tab:methods} contrasts the surveyed methods along the two
challenges identified above. No single method addresses both camouflage
and imbalance in a simple, stable way: FRAUDAR handles camouflage but
ignores features and labels; CARE-GNN handles camouflage with an
expensive RL selector; PC-GNN handles imbalance at the sampling level
but ignores camouflage. Our approach (Section~IV) combines a fixed,
label-supervised similarity filter in the spirit of CARE-GNN --- removing
the RL machinery --- with PC-GNN-style balanced neighborhood sampling and
a class-weighted objective, and we ablate each component to measure its
contribution.

\begin{table}[t]
\centering
\caption{Surveyed methods and the challenges they address.}
\label{tab:methods}
\begin{tabular}{llcc}
\hline
\textbf{Method} & \textbf{Venue} & \textbf{Camouflage} & \textbf{Imbalance} \\
\hline
FRAUDAR~\cite{DBLP:conf/kdd/HooiSBSSF16} & KDD'16 & yes & no \\
SemiGNN~\cite{DBLP:conf/icdm/WangQL0JWFYZY19} & ICDM'19 & no & partial \\
CARE-GNN~\cite{DBLP:conf/cikm/DouL0DPY20} & CIKM'20 & yes & no \\
PC-GNN~\cite{DBLP:conf/www/LiuAQCFYH21} & WWW'21 & no & yes \\
\textbf{Ours} & -- & yes & yes \\
\hline
\end{tabular}
\end{table}

\section{Problem Definition}
Let $G = (\mathcal{V}, \mathcal{E}, \mathbf{X})$ be an attributed graph with
node set $\mathcal{V}$, edge set $\mathcal{E}$, and node feature matrix
$\mathbf{X} \in \mathbb{R}^{|\mathcal{V}| \times d}$. In our benchmarks the
graph is multi-relation: $\mathcal{E}$ is the union of relation-specific
edge sets $\mathcal{E}_1, \dots, \mathcal{E}_R$ ($R = 3$). Each node
$v \in \mathcal{V}$ has a label $y_v \in \{0, 1\}$ ($1$ = fraudster), of which
only a small labeled subset $\mathcal{V}_L \subset \mathcal{V}$ is observed,
with $|\{v : y_v = 1\}| \ll |\{v : y_v = 0\}|$ (class imbalance). The task is
\emph{semi-supervised node classification}: learn
$f : \mathcal{V} \rightarrow [0, 1]$ from $G$ and $\mathcal{V}_L$ that ranks
fraudulent nodes above benign ones on the unlabeled nodes.

\section{Proposed Method}
Our model, an \emph{imbalance-aware multi-relation GNN}, has three
components: (A) camouflage-resistant neighbor filtering, (B) balanced
neighborhood sampling, and (C) relation-wise attention aggregation and
fusion. We describe each component formally.

\subsection{Camouflage-Resistant Neighbor Filtering}
Following CARE-GNN~\cite{DBLP:conf/cikm/DouL0DPY20}, we assume camouflaged
neighbors are feature-wise dissimilar to the center node. For each relation
$r$, a two-layer MLP $\phi_r$ scores the similarity between the center node
$u$ and a neighbor $v$ from their projected features
$\mathbf{h}_u, \mathbf{h}_v$:
\begin{equation}
s_r(u, v) = \sigma\!\left(\phi_r\big([\mathbf{h}_u \,\|\, \mathbf{h}_v]\big)\right).
\end{equation}
Instead of CARE-GNN's RL controller, we keep a fixed fraction $\rho = 0.5$
of each node's neighbors per relation (the top-scored ones), which removes
the RL training loop entirely and is stable in practice. Because hard top-$k$
selection is non-differentiable, $\phi_r$ receives no gradient from the
classification loss; as in CARE-GNN, we therefore train it with label
supervision: a binary cross-entropy loss $\mathcal{L}_{\text{sim}}$ that
predicts, for edges incident to labeled training nodes, whether the two
endpoints share the same label.

\subsection{Balanced Neighborhood Sampling}
Following PC-GNN~\cite{DBLP:conf/www/LiuAQCFYH21}, we rebalance the
aggregated neighborhood toward the minority class. Each node's filtered
neighborhood is capped at $K = 32$ neighbors, sampled without replacement
with probability proportional to a weight $w_v$: $w_v = \beta = 4$ for
neighbors known (from \emph{training} labels) to be fraudulent, and $w_v = 1$
otherwise. Sampling uses Efraimidis--Spirakis keys $\log u / w_v$,
$u \sim \mathcal{U}(0,1)$, which gives exact weighted sampling without
replacement; at evaluation time sampling is deterministic (key $\log w_v$).

\subsection{Relation-wise Aggregation and Fusion}
Within each relation, the filtered and sampled neighborhood is aggregated
with dot-product attention:
\begin{equation}
\alpha_{uv} = \mathrm{softmax}_{v}\!\left(\frac{\mathbf{h}_u^\top
\mathbf{h}_v}{\sqrt{d}}\right), \qquad
\mathbf{h}_u^{(r)} = \sum_{v \in \hat{\mathcal{N}}_r(u)} \alpha_{uv}\,
\mathbf{W} \mathbf{h}_v,
\end{equation}
where $\hat{\mathcal{N}}_r(u)$ is the filtered, sampled neighborhood. A
residual connection preserves the center node's own information,
$\tilde{\mathbf{h}}_u^{(r)} = \mathbf{h}_u^{(r)} + \mathbf{h}_u$, and a
learned cross-relation attention vector fuses the $R$ relation embeddings
into the final representation, which a linear layer maps to two logits.

\subsection{Training Objective}
The model is trained with class-weighted cross-entropy
$\mathcal{L}_{\text{CE}}^{w}$ (inverse class frequency, computed on the
training split) plus the similarity-supervision term:
\begin{equation}
\mathcal{L} = \mathcal{L}_{\text{CE}}^{w} + \lambda \,
\mathcal{L}_{\text{sim}}, \qquad \lambda = 0.1.
\end{equation}
All components are trained end-to-end with Adam.

\section{Experiments}
\subsection{Datasets}
We use two public multi-relation fraud benchmarks released with
CARE-GNN~\cite{DBLP:conf/cikm/DouL0DPY20} (Table~\ref{tab:datasets}).
\textbf{YelpChi} contains hotel and restaurant reviews labeled as spam or
genuine, with relations R-U-R (same user), R-S-R (same star rating), and
R-T-R (same product and time window). \textbf{Amazon} contains product
reviewers labeled as fraudsters or benign, with relations U-P-U
(co-reviewed product), U-S-U (same star rating), and U-V-V (similar review
text). Following CARE-GNN, features are L2-normalized, splits are stratified
with 40\% train, and Amazon's unlabeled nodes are excluded from the splits
but kept in the graph for message passing; we carve the remaining 60\% into
20\% validation and 40\% test.

\begin{table}[t]
\centering
\caption{Dataset statistics (fraud rate is over labeled nodes).}
\label{tab:datasets}
\begin{tabular}{lcc}
\hline
 & \textbf{YelpChi} & \textbf{Amazon} \\
\hline
Nodes & 45{,}954 & 11{,}944 \\
Edges (all relations) & 7{,}693{,}758 & 8{,}796{,}784 \\
Features & 32 & 25 \\
Fraud rate & 14.5\% & 9.5\% \\
\hline
\end{tabular}
\end{table}

\subsection{Experimental Setup}
We compare our model against two baselines: a two-layer \textbf{GCN} on the
union graph, and \textbf{SemiGNN}~\cite{DBLP:conf/icdm/WangQL0JWFYZY19}
(implemented per the paper, with its semi-supervised auxiliary loss). All
models use hidden dimension 64, Adam with learning rate 0.01, 100 epochs,
and class-weighted cross-entropy; the best checkpoint by validation AUROC
is evaluated on the test split. We report AUROC, AUPRC, F1-macro, and
Recall@100, the standard metrics for this
benchmark~\cite{DBLP:conf/cikm/DouL0DPY20,DBLP:conf/www/LiuAQCFYH21}, as
mean $\pm$ standard deviation over three seeds (1, 7, 42). All experiments
run on CPU; a full training run takes 15--60 minutes depending on the model.

\subsection{Main Results}
Table~\ref{tab:main} shows the main results. Both relation-aware models
vastly outperform GCN on YelpChi (GCN remains close to random at 0.585
AUROC), confirming that flat homophily-based aggregation fails under
camouflage. On YelpChi our model is the best on every metric: it improves
over SemiGNN by $+5.3$ AUROC and $+15.2$ F1-macro, a gap far larger than
the seed-to-seed standard deviation ($\leq 0.01$). On Amazon, SemiGNN keeps
a narrow edge in ranking metrics (AUROC/AUPRC), while our model achieves the
best F1-macro and Recall@100. We attribute the difference to the structure
of the Amazon graph, where a single relation (U-S-U) carries over 80\% of
edges, leaving less room for similarity filtering to help.

\begin{table}[t]
\centering
\caption{Test results (mean $\pm$ std over 3 seeds). Best per column in bold.}
\label{tab:main}
\begin{tabular}{llccc}
\hline
\textbf{Dataset} & \textbf{Model} & \textbf{AUROC} & \textbf{AUPRC} & \textbf{F1-macro} \\
\hline
\multirow{3}{*}{YelpChi}
 & GCN      & 0.585$\pm$.002 & 0.216$\pm$.008 & 0.443$\pm$.088 \\
 & SemiGNN  & 0.801$\pm$.008 & 0.407$\pm$.012 & 0.540$\pm$.000 \\
 & Ours     & \textbf{0.854$\pm$.003} & \textbf{0.518$\pm$.006} & \textbf{0.692$\pm$.003} \\
\hline
\multirow{3}{*}{Amazon}
 & GCN      & 0.855$\pm$.002 & 0.444$\pm$.008 & 0.621$\pm$.012 \\
 & SemiGNN  & \textbf{0.959$\pm$.004} & \textbf{0.879$\pm$.009} & 0.878$\pm$.011 \\
 & Ours     & 0.944$\pm$.002 & 0.858$\pm$.004 & \textbf{0.904$\pm$.003} \\
\hline
\end{tabular}
\end{table}

\subsection{Ablation Study}
Table~\ref{tab:ablation} ablates the two components of our model.
Removing the similarity filter (Ours$_{\neg\text{filter}}$, which also drops
$\sim$25K parameters) degrades YelpChi AUROC by $4.5$ points and F1-macro by
$8.1$ points, showing that camouflage-resistant filtering is the main driver
of the improvement. On Amazon, removing the filter slightly raises
AUROC but cuts F1-macro by $12$ points: filtering mainly improves
\emph{decision} quality rather than \emph{ranking} quality there. Removing
balanced sampling (Ours$_{\neg\text{sampler}}$) changes results by less than
$0.1$ AUROC points on both datasets --- the component contributes nothing in
this setting, a negative result we analyze next.

\begin{table}[t]
\centering
\caption{Ablation of our model (seed 42).}
\label{tab:ablation}
\begin{tabular}{llccc}
\hline
\textbf{Dataset} & \textbf{Variant} & \textbf{AUROC} & \textbf{AUPRC} & \textbf{F1-macro} \\
\hline
\multirow{3}{*}{YelpChi}
 & Ours (full)              & 0.856 & 0.521 & 0.694 \\
 & Ours$_{\neg\text{filter}}$  & 0.811 & 0.449 & 0.613 \\
 & Ours$_{\neg\text{sampler}}$ & 0.855 & 0.519 & 0.699 \\
\hline
\multirow{3}{*}{Amazon}
 & Ours (full)              & 0.943 & 0.853 & 0.902 \\
 & Ours$_{\neg\text{filter}}$  & 0.948 & 0.856 & 0.779 \\
 & Ours$_{\neg\text{sampler}}$ & 0.943 & 0.853 & 0.902 \\
\hline
\end{tabular}
\end{table}

\subsection{Interaction of Imbalance Mechanisms}
To test whether balanced sampling is merely redundant with class-weighted
training, we ran a $2 \times 2$ study crossing the two imbalance mechanisms
(sampling vs.\ class-weighted loss). With class weights disabled, the
sampler still changes results by $\leq 0.2$ AUROC points on both datasets,
while removing class weights on YelpChi collapses F1-macro from 0.69 to 0.46.
We conclude that (i) loss-level reweighting is the effective imbalance
mechanism on these benchmarks, and (ii) PC-GNN-style sampling does not
transfer to this setting, plausibly because known-fraud neighbors are too
rare (9--15\% base rate) for a $4\times$ sampling boost to alter aggregated
neighborhoods materially. We view reporting such negative results as
important for reproducible fraud-detection research.

\section{Conclusion and Future Work}
We surveyed representative graph-based fraud detection methods and proposed
a simple imbalance-aware GNN that replaces CARE-GNN's RL neighbor selector
with a fixed, label-supervised similarity filter. The model is the best
performer on YelpChi across all metrics and achieves the best F1-macro on
Amazon, and our studies identify the similarity filter as the decisive
component. Future work includes heterophily-aware aggregation, tuning or
removing the balanced sampler (e.g., larger boost factors or learnable
sampling), extending the ablations to multiple seeds, and evaluating on
additional datasets such as Elliptic.

\bibliographystyle{IEEEtran}
\bibliography{references}

\end{document}
